\section{引言与主要结论}
\label{sec:introduction}

实欧氏空间中的 Kakeya 集合，是指在每个无向方向上都包含一条单位线段的集合。
Kakeya 猜想预言，这种集合在 $\R^n$ 中的 Hausdorff 维数为 $n$。
本文研究的对象始终是实欧氏空间中普通的直线 Kakeya 集合，
不把有限域版本、Minkowski 维数、受限方向结论或附加结构的集合类与之混同。

四维原稿\cite{OAI4}声称四维 Kakeya 集合具有满 Hausdorff 维数。
对高维转移而言，仅有这一集合维数声明尚不足够：递推需要一项能对任意有限线族、
原着色及任意固定环境维数统一调用的定量输入（常数可依赖该环境维数）。
本文使用的是该证明的有限尺度读出及其稀疏方向放大，明确表述为引用输入\ref{input:D4}。
其证明依赖与档案定位在附录\ref{app:inputs}说明。
既有有限证明和成果报告分别见\cite{ArchiveFSD,ArchiveDIM}。

\begin{theorem}[六维下界]\label{thm:main}
采用四维有限输入\ref{input:D4}和 Katz--Tao 递推定理\ref{input:KT}。
若任意集合 $K\subset\R^6$ 在每个欧氏方向上包含一条单位线段，则
\begin{equation}\label{eq:main}
  \HH K\ge\frac{17}{4}=4.25.
\end{equation}
不要求 $K$ 可测、紧或有界，不要求存在可测的方向选线映射。
\end{theorem}

定理的证明分为三个步骤。第一，利用随机线性投影把任意固定 $\R^M$ 中的有限线族送到 $\R^4$。
投影方向可能发生碰撞，投影线段也可能变短；两类损失都必须在使用四维输入之前支付。
第二，保留任意原条数的非饱和估计，得到每个固定实参数 $n\ge4$ 的严格次临界种子
$\Kest(n,4-\zeta)$。对 $0<\zeta<1$，Katz--Tao 定理给出
\begin{equation}\label{eq:roadmap-recurrence}
  \Kest(6,4-\zeta),\quad\Kest(5-\zeta,4-\zeta)
  \quad\Longrightarrow\quad
  \Kest\!\left(6,\frac{17}{4}-\frac{\zeta}{4}\right).
\end{equation}
第三，以同尺度覆盖和外方向测度建立任意集合的 Hausdorff 读出，再令 $\zeta\downarrow0$。

Katz--Tao\cite{KT}的经典六维下界为 $7-2\sqrt2$。
Hickman--Rogers--Zhang\cite[Figure 5]{HRZ}以及 Gao--Liu--Xi\cite[Table 1]{GLX}
的同类型欧氏对照亦明确列出该数值。二者之差为
\begin{equation}\label{eq:comparison}
 \frac{17}{4}-(7-2\sqrt2)
   =2\sqrt2-\frac{11}{4}
   \approx0.078427125>0.
\end{equation}
这一比较只针对上述明确核到的普通六维 Hausdorff 基线，不是对全部当前文献的最优性认证。
本文的贡献归属是四维有限输入的高维后果及所需有限接口的整理；递推本身归于 Katz--Tao。

\paragraph{全文结构。}
第\ref{sec:preliminaries}节固定离散约定与两个具名输入。
第\ref{sec:projection}节证明有限随机投影和原着色保留。
第\ref{sec:seeds}节给出非饱和估计及真实递推参数。
第\ref{sec:hausdorff}节证明任意集合的 Hausdorff 读出并完成定理\ref{thm:main}。
第\ref{sec:scope}节说明范围；附录列出四维输入来源及稀疏化约。
